\documentclass[a4paper, 11pt]{article}

\usepackage[ngerman]{babel}
\usepackage[margin=2.5cm]{geometry}
\usepackage{amsopn}

\title{Implementierungsanleitung}
\author{}
\date{}

\begin{document}

    \maketitle

    \section{Ziehle dieses Algorithmus}
    Diese Implementation soll mit möglichst wenig Speicher auskommen, aber trotzdem eine gute Laufzeit haben.

% Text

    \section{Implementierung des Algorithmus}
    Der Algorithmus erhält als Übergabeparameter zwei Arrays: \textit{Numbers} und \textit{Targets}.\\\\
    Zunächst muss der \textit{Numbers} Array sortiert werden, dies wird für die spätere Binärsuche benötigt: Hier bieten sich Suchalgorithmen wie Quicksort o. Ä an.\\\\
    Nun beginnt der Hauptteil des Algorithmus: Für jede Zahl \textit{target} in \textit{Targets} wird nun folgendermaßen vorgegangen: Für jede Zahl \textit{number} in \textit{Numbers} wird die Differenz \textit{target} - \textit{number} gebildet und zwischengespeichert.\\
    Um jetzt herauszufinden, ob die Zielzahl \textit{target} aus zwei Zahlen in \textit{Numbers} gebildet werden kann, muss man schauen, ob sich diese Differenz in \textit{Numbers} befindet. Hierzu bietet sich die Binärsuche an, die die gesuchte Zahl (falls vorhanden) in \textit{Numbers} findet.
% Text

    \section{Weitere Optimierungen}


% Text

    \section{Laufzeit}
    Für alle n Zahlen in \textit{Targets} durchlaufen wir alle m Zahlen in \textit{Numbers}. Für jede Zahl m muss nun final das Gegenstück in \textit{Numbers} gesucht werden, was mit der Binärsuche eine Laufzeit von log(m) hat.\\\\
    Zusammen: \[
                  \mathcal{O}\left(n \cdot m \cdot \log(m)\right)
    \]

% Text

\end{document}